\documentclass[11pt, a4paper]{article}

% bfw-style.tex — gemeinsame Style-Definitionen für alle Handouts/Aufgabenhefte
% Voraussetzung: Kompilation mit xelatex (wegen fontspec)

\usepackage[ngerman]{babel}
% Babel-ngerman macht " zu einem aktiven Shorthand (z.B. für "a -> ä). In unseren
% Texten verwenden wir " als normales Zeichen — daher Shorthand komplett ausschalten.
\shorthandoff{"}
\usepackage{geometry}
\geometry{a4paper, top=2.4cm, bottom=2.4cm, left=2.3cm, right=2.3cm,
          headheight=15pt, headsep=0.7cm, footskip=1.1cm}

\usepackage{fontspec}
% Fallback-Kette: Source Sans Pro -> Calibri -> Segoe UI -> Latin Modern Sans
% (Latin Modern ist immer mit MiKTeX vorhanden)
\IfFontExistsTF{Source Sans Pro}{%
  \setmainfont{Source Sans Pro}[Ligatures=TeX]%
}{\IfFontExistsTF{Calibri}{%
  \setmainfont{Calibri}[Ligatures=TeX]%
}{\IfFontExistsTF{Segoe UI}{%
  \setmainfont{Segoe UI}[Ligatures=TeX]%
}{%
  \setmainfont{Latin Modern Sans}[Ligatures=TeX]%
}}}
\IfFontExistsTF{Source Code Pro}{%
  \setmonofont{Source Code Pro}[Scale=0.92]%
}{\IfFontExistsTF{Consolas}{%
  \setmonofont{Consolas}[Scale=0.92]%
}{%
  \setmonofont{Latin Modern Mono}[Scale=0.92]%
}}

% Gesperrte Versalien für Labels/Titelbalken (Kapitälchen-Ersatz, funktioniert
% mit jeder OpenType-Schrift — Latin Modern Sans hat keine echten Kapitälchen)
\newcommand{\bfwlabelfont}{\footnotesize\bfseries\addfontfeatures{LetterSpace=8.0}}

\usepackage{xcolor}
% Farbpalette: hell, freundlich, modern
\definecolor{bfwPrimary}{HTML}{0EA5A4}    % Petrol/Türkis - Primär
\definecolor{bfwPrimaryDark}{HTML}{0F766E} % dunkler Petrol für Headings
\definecolor{bfwAccent}{HTML}{F59E0B}     % Amber - Akzent
\definecolor{bfwSuccess}{HTML}{10B981}    % Grün - Erfolg / Lösung
\definecolor{bfwWarn}{HTML}{F97316}       % Orange - Achtung
\definecolor{bfwInk}{HTML}{0F172A}        % fast Schwarz
\definecolor{bfwInkSoft}{HTML}{475569}    % gedämpftes Grau
\definecolor{bfwBgSoft}{HTML}{F8FAFC}     % off-white
\definecolor{bfwBgTeal}{HTML}{ECFEFF}     % helles Türkis
\definecolor{bfwBgAmber}{HTML}{FEF3C7}    % helles Gelb
\definecolor{bfwBgRose}{HTML}{FFE4E6}     % helles Rosé für Eselsbrücken
\definecolor{bfwTabHead}{HTML}{E4F3F2}    % zarte Kopfzeilen-Hinterlegung für Tabellen

\usepackage[colorlinks=true, linkcolor=bfwPrimaryDark, urlcolor=bfwPrimaryDark]{hyperref}
\usepackage{lastpage}
\usepackage{enumitem}
\setlist[itemize]{leftmargin=*, itemsep=2pt, topsep=2pt}
\setlist[enumerate]{leftmargin=*, itemsep=2pt, topsep=2pt}

\usepackage[most]{tcolorbox}
\usepackage{booktabs}
\usepackage{colortbl}
\usepackage{tikz}
\usetikzlibrary{decorations.pathmorphing, positioning, shapes.geometric, arrows.meta}

% ---------------------------------------------------------------------------
% Einheitliches Tabellen-Design
%   - etwas mehr Zeilenluft, dezente graue Linien (wirkt auch mit \hline ruhiger)
%   - Kopfzeile einer Tabelle bei Bedarf zart hinterlegen: \rowcolor{bfwTabHead}
% ---------------------------------------------------------------------------
\renewcommand{\arraystretch}{1.25}
\arrayrulecolor{bfwInkSoft!50}
\setlength{\heavyrulewidth}{1pt}
\setlength{\lightrulewidth}{0.5pt}

% ---------------------------------------------------------------------------
% Boxen — klare farbige Titelbalken mit gesperrten Versal-Labels
% (Titel bleibt per Option überschreibbar: \begin{merksatz}[title={...}])
% ---------------------------------------------------------------------------
\tcbset{bfwbox/.style={%
  enhanced, breakable,
  boxrule=0.5pt, arc=2.5pt,
  left=10pt, right=10pt, top=8pt, bottom=8pt,
  toptitle=4pt, bottomtitle=4pt,
  titlerule=0pt,
  fonttitle=\bfwlabelfont,
  coltitle=white
}}

% Merkkästchen — wichtige Definition / Kernpunkt
\newtcolorbox{merksatz}[1][]{%
  bfwbox,
  colback=bfwBgTeal,
  colframe=bfwPrimaryDark,
  colbacktitle=bfwPrimaryDark,
  title={MERKE},
  #1
}

% Eselsbrücke — Mnemoniks
\newtcolorbox{eselsbruecke}[1][]{%
  bfwbox,
  colback=bfwBgRose,
  colframe=bfwAccent!85!black,
  colbacktitle=bfwAccent!85!black,
  title={ESELSBRÜCKE},
  #1
}

% Beispiel-Box
\newtcolorbox{beispiel}[1][]{%
  bfwbox,
  colback=bfwBgSoft,
  colframe=bfwInkSoft,
  colbacktitle=bfwInkSoft,
  title={BEISPIEL},
  #1
}

% Lösungs-Box (für Aufgabenhefte)
\newtcolorbox{loesung}[1][]{%
  bfwbox,
  colback=bfwSuccess!7,
  colframe=bfwSuccess!65!black,
  colbacktitle=bfwSuccess!65!black,
  title={LÖSUNG},
  #1
}

% Achtung-Box — typische Fehler / Stolperfallen
\newtcolorbox{achtung}[1][]{%
  bfwbox,
  colback=bfwWarn!6,
  colframe=bfwWarn!85!black,
  colbacktitle=bfwWarn!85!black,
  title={ACHTUNG},
  #1
}

% Formel-Box — für zentrale Formeln (groß, ruhig, ohne Titelbalken)
\newtcolorbox{formelbox}[1][]{%
  enhanced,
  colback=bfwBgTeal,
  colframe=bfwPrimaryDark,
  boxrule=1pt, arc=3pt,
  left=10pt, right=10pt, top=6pt, bottom=6pt,
  #1
}

% Glossar-Eintrag
\newcommand{\glossareintrag}[2]{%
  \par\noindent\textbf{\color{bfwPrimaryDark}#1} \enspace #2\par\smallskip
}

% ---------------------------------------------------------------------------
% Abschnitts-Design: farbiges Nummern-Quadrat + feine Linie unter \section,
% dezent farbige \subsection/\subsubsection
% ---------------------------------------------------------------------------
\usepackage{titlesec}
\newcommand{\bfwSecBadge}[1]{%
  \begin{tikzpicture}[baseline={([yshift=-0.62em]n.center)}]
    \node[fill=bfwPrimaryDark, text=white, font=\normalsize\bfseries,
          minimum width=1.55em, minimum height=1.55em, inner sep=1pt,
          rounded corners=1.5pt] (n) {#1};
  \end{tikzpicture}}
\titleformat{\section}
  {\Large\bfseries\color{bfwPrimaryDark}}
  {\bfwSecBadge{\thesection}}{0.6em}{}
  [\vspace{-0.2em}{\color{bfwPrimary!60}\titlerule[0.8pt]}]
\titleformat{\subsection}
  {\large\bfseries\color{bfwPrimaryDark}}
  {\textcolor{bfwPrimary}{\thesubsection}}{0.5em}{}
\titleformat{\subsubsection}
  {\normalsize\bfseries\color{bfwInk}}
  {\textcolor{bfwPrimary}{\thesubsubsection}}{0.4em}{}
\titlespacing*{\section}{0pt}{1.6em}{1em}
\titlespacing*{\subsection}{0pt}{1.2em}{0.5em}

% TikZ-Stil "skizzenhaft" — etwas weicher als Rough.js, aber stabil
\tikzset{
  roughbox/.style = {
    draw=bfwPrimaryDark, thick, fill=bfwBgTeal,
    rounded corners=4pt, inner sep=6pt
  },
  roughaccent/.style = {
    draw=bfwAccent, thick, fill=bfwBgAmber,
    rounded corners=4pt, inner sep=6pt
  },
  roughline/.style = {
    draw=bfwInkSoft, thick, -{Stealth[length=2.5mm]}
  }
}

% Bullet-Symbole (bewusst kein FontAwesome — vermeidet Font-Installations-Probleme)
\newcommand{\faLightbulb}{$\bullet$}
\newcommand{\faBookmark}{$\bullet$}

% ---------------------------------------------------------------------------
% Kopf-/Fußzeile
%   Kopf links:  Wortlogo „Wirtschaftsfachwirt"
%   Kopf rechts: Dokument-Kurztext, gesetzt via \kopfzeile{...}
%   Fuß links:   © Can Siebert · 2026 — Fuß rechts: Seite X von Y
%   Titelseite (Seite 1) bleibt ohne Kopf-/Fußzeile (\handouttitel setzt
%   \thispagestyle{empty}).
% ---------------------------------------------------------------------------
\usepackage{fancyhdr}
\newcommand{\bfwKopfzeileText}{}
\newcommand{\kopfzeile}[1]{\renewcommand{\bfwKopfzeileText}{#1}}
\pagestyle{fancy}
\fancyhf{}
\fancyhead[L]{\small\bfseries\color{bfwPrimaryDark}Wirtschaftsfachwirt}
\fancyhead[R]{\small\color{bfwInkSoft}\bfwKopfzeileText}
\renewcommand{\headrulewidth}{0.6pt}
\renewcommand{\headrule}{{\color{bfwPrimary}\hrule width\headwidth height 0.6pt}}
\fancyfoot[L]{\small\color{bfwInkSoft}\copyright\ Can Siebert\,\textperiodcentered\,2026}
\fancyfoot[R]{\small\color{bfwInkSoft}Seite \thepage\ von \pageref*{LastPage}}
% Auch \thispagestyle{plain} (z.B. durch \maketitle) soll leer bleiben:
\fancypagestyle{plain}{\fancyhf{}\renewcommand{\headrulewidth}{0pt}\renewcommand{\headrule}{}}

% ---------------------------------------------------------------------------
% \handouttitel{Obertitel}{Untertitel}{Kurs-Label}{Termin-Zeile}
%   Einheitlicher professioneller Titelblock: Dachzeile, farbige Headline,
%   Untertitel, farbige Trennlinie und dezente Meta-Zeile (Kurs/Termin/Dozent).
%   Ersetzt \title/\maketitle. Direkt nach \begin{document} aufrufen.
% ---------------------------------------------------------------------------
\newcommand{\handouttitel}[4]{%
  \thispagestyle{empty}%
  \begingroup
  \setlength{\parindent}{0pt}%
  {\bfwlabelfont\color{bfwPrimary}WIRTSCHAFTSFACHWIRT (IHK)\par}
  \vspace{0.55em}
  {\huge\bfseries\color{bfwPrimaryDark}#1\par}
  \vspace{0.5em}
  {\large\color{bfwInkSoft}#2\par}
  \vspace{0.9em}
  {\color{bfwPrimary}\rule{\linewidth}{2pt}}\par
  \vspace{0.55em}
  {\small\color{bfwInkSoft}#3\ \,\textperiodcentered\,\ #4\ \,\textperiodcentered\,\ Dozent: Can Siebert\par}
  \endgroup
  \vspace{1.4em}%
}


\kopfzeile{Aufgabenheft BWL Lektion 1}

\begin{document}
\handouttitel{Aufgabenheft Lektion~1 (BWL)}%
  {Betriebliche Funktionen I: Produktion, Logistik \& Marketing --- Übungen im IHK-Klausurformat}%
  {1.2 Betriebswirtschaftliche Grundlagen}%
  {Selbstlernmaterial zur Lektion vom Sa, 05.09.2026}


\begin{merksatz}[title={\bfseries So nutzen Sie dieses Aufgabenheft}]
Bearbeiten Sie alle Aufgaben zunächst \emph{ohne} in die Lösungen zu schauen. Schreiben
Sie Ihre Antworten in vollständigen Sätzen --- die IHK-Prüfung verlangt formulierte
Begründungen, nicht bloße Stichworte. Beachten Sie die \emph{Operatoren} (nennen,
erläutern, berechnen, begründen, entscheiden) --- sie geben den geforderten
Antwort-Tiefegrad vor. Erst danach öffnen Sie den Lösungsteil ab
Seite~\pageref{loesungen} und vergleichen. Ergänzend: Übungsband Betriebswirtschaft,
Kapitel~1 (Übungen 1--18).
\end{merksatz}

\tableofcontents
\vspace{1em}
\hrule

\section{Aufgaben}

\subsection*{Aufgabe~1 --- Zielsystem des Unternehmens (10 Punkte)}

Die Geschäftsleitung der Nordwest Küchengeräte AG formuliert für das kommende Jahr:
„Wir wollen hochwertige Küchengeräte fertigen (1), die Umsatzrendite auf 8~\% steigern
(2), die Ausschussrate halbieren (3) und den CO$_2$-Ausstoß der Fertigung um 20~\%
senken (4).``

\begin{enumerate}[label=(\alph*)]
  \item \textbf{Unterscheiden} Sie Sachziele und Formalziele und \textbf{ordnen} Sie die Aussagen (1) und (2) zu. \hfill (4~P.)
  \item \textbf{Nennen} Sie die drei möglichen Zielbeziehungen zwischen Unternehmenszielen. \hfill (3~P.)
  \item \textbf{Beurteilen} Sie die Zielbeziehung zwischen Ziel (3) und Ziel (2) sowie zwischen Ziel (4) und Ziel (2) --- jeweils mit kurzer Begründung. \hfill (3~P.)
\end{enumerate}

\subsection*{Aufgabe~2 --- Produktionsfaktoren nach Gutenberg (12 Punkte)}

\begin{enumerate}[label=(\alph*)]
  \item \textbf{Grenzen} Sie die betriebswirtschaftlichen Produktionsfaktoren nach Gutenberg von den volkswirtschaftlichen Produktionsfaktoren \textbf{ab}. \hfill (4~P.)
  \item \textbf{Nennen} Sie die drei Elementarfaktoren und die Aufgaben des dispositiven Faktors. \hfill (4~P.)
  \item In einem Möbelwerk fallen an: Holz, Leim, Schmieröl für die Sägen, eine CNC-Fräse. \textbf{Ordnen} Sie die vier Positionen den Faktoren Rohstoff, Hilfsstoff, Betriebsstoff und Betriebsmittel zu. \hfill (4~P.)
\end{enumerate}

\subsection*{Aufgabe~3 --- Fertigungstypen (10 Punkte)}

Ordnen Sie die folgenden Situationen jeweils dem zutreffenden Fertigungstyp
(Einzel-, Serien-, Sorten-, Massenfertigung) zu und \textbf{begründen} Sie kurz:

\begin{enumerate}[label=(\alph*)]
  \item Eine Werft baut ein Kreuzfahrtschiff nach individuellem Kundenauftrag. \hfill (2~P.)
  \item Eine Brauerei stellt auf derselben Anlage nacheinander Pils, Export und Weizenbier her. \hfill (2~P.)
  \item Ein Automobilhersteller fertigt von einem Modell 30.000 Stück, rüstet dann auf ein anderes Modell um. \hfill (2~P.)
  \item Ein Schraubenhersteller produziert eine Standardschraube ununterbrochen in Millionenauflage. \hfill (2~P.)
  \item \textbf{Unterscheiden} Sie Serien- und Sortenfertigung anhand des entscheidenden Merkmals. \hfill (2~P.)
\end{enumerate}

\subsection*{Aufgabe~4 --- Produktivität und Wirtschaftlichkeit (14 Punkte)}

Die Nordwest Küchengeräte AG fertigte im Monat August mit 20 Mitarbeitern
(je 160 Arbeitsstunden) 8.000 Mixergehäuse. Der bewertete Ertrag betrug 400.000~€,
der Aufwand 320.000~€. Im September wurden mit gleicher Stundenzahl 8.800 Gehäuse
gefertigt; der Ertrag stieg auf 440.000~€, der Aufwand auf 352.000~€.

\begin{enumerate}[label=(\alph*)]
  \item \textbf{Erklären} Sie den Unterschied zwischen Produktivität und Wirtschaftlichkeit. \hfill (4~P.)
  \item \textbf{Berechnen} Sie die Arbeitsproduktivität (Stück je Arbeitsstunde) für August und September. \hfill (4~P.)
  \item \textbf{Berechnen} Sie die Wirtschaftlichkeit für beide Monate und \textbf{interpretieren} Sie die Ergebnisse. \hfill (4~P.)
  \item \textbf{Nennen} Sie das dritte zentrale Produktionsziel neben Wirtschaftlichkeit und Produktivität. \hfill (2~P.)
\end{enumerate}

\subsection*{Aufgabe~5 --- Logistik: Teilbereiche und 6 R (12 Punkte)}

\begin{enumerate}[label=(\alph*)]
  \item \textbf{Nennen} Sie die 6 R der Logistik. \hfill (6~P.)
  \item Die folgende Beschreibung enthält sachliche Fehler: „Die Beschaffungslogistik beschäftigt sich mit dem Informations- und Materialfluss vom Kunden zum Unternehmen. Sie betrifft die Abläufe im Lager und in der Produktion.`` \textbf{Stellen} Sie die Aussage \textbf{richtig}. \hfill (3~P.)
  \item \textbf{Bringen} Sie den Güterdurchlauf in die richtige Reihenfolge: Absatzlager --- Produktion --- Versand --- Materiallager (vom Lieferanten zum Kunden). \hfill (3~P.)
\end{enumerate}

\subsection*{Aufgabe~6 --- Lagerfunktionen und Supply Chain Management (12 Punkte)}

\begin{enumerate}[label=(\alph*)]
  \item \textbf{Nennen und erläutern} Sie vier Funktionen der Lagerhaltung mit je einem Beispiel. \hfill (8~P.)
  \item \textbf{Beschreiben} Sie, was unter Supply Chain Management zu verstehen ist, und \textbf{nennen} Sie zwei Ziele. \hfill (4~P.)
\end{enumerate}

\subsection*{Aufgabe~7 --- Marketingbegriff, Marktforschung und Marketingziele (12 Punkte)}

\begin{enumerate}[label=(\alph*)]
  \item \textbf{Beurteilen} Sie die Aussage: „Ausgangspunkt für alle Marketingaktivitäten ist das Produkt.`` \hfill (3~P.)
  \item \textbf{Unterscheiden} Sie Primär- und Sekundärforschung und \textbf{nennen} Sie je zwei Methoden bzw. Quellen. \hfill (5~P.)
  \item \textbf{Charakterisieren} Sie ökonomische und psychografische Marketingziele und \textbf{nennen} Sie je zwei Beispiele. \hfill (4~P.)
\end{enumerate}

\subsection*{Aufgabe~8 --- Fallaufgabe Marketing-Mix (16 Punkte)}

Die Nordwest Küchengeräte AG bringt einen neuen Hochleistungsmixer „VitaMix Pro`` auf
den Markt. Zielgruppe sind gesundheitsbewusste Haushalte mit mittlerem bis hohem
Einkommen. Der Vertriebsleiter schlägt vor, den Mixer ausschließlich über den eigenen
Onlineshop zu verkaufen; die Marketingabteilung plant eine Einführungskampagne.

\begin{enumerate}[label=(\alph*)]
  \item \textbf{Nennen} Sie die vier Instrumentalbereiche des Marketing-Mix (deutsch, ggf. mit englischem Fachbegriff). \hfill (4~P.)
  \item \textbf{Entwickeln} Sie für den „VitaMix Pro`` zu jedem der vier Instrumente je zwei konkrete Maßnahmen bzw. Entscheidungen. \hfill (8~P.)
  \item Der Mixer soll „zum richtigen Zeitpunkt in der gewünschten Menge am Ort der Nachfrage`` verfügbar sein. \textbf{Ordnen} Sie diese Aufgabe dem zutreffenden Instrument zu und \textbf{begründen} Sie. \hfill (2~P.)
  \item \textbf{Unterscheiden} Sie Werbung und Public Relations. \hfill (2~P.)
\end{enumerate}

\subsection*{Aufgabe~9 --- Erfolgsrechnung zur Produkteinführung (10 Punkte)}

Für den „VitaMix Pro`` sind bekannt: Fixkosten 150.000~€, variable Kosten 14~€/Stück,
geplanter Marktpreis 30~€/Stück, erwartete Absatzmenge 150.000 Stück.

\begin{enumerate}[label=(\alph*)]
  \item \textbf{Berechnen} Sie den erwarteten Gewinn. \hfill (4~P.)
  \item \textbf{Berechnen} Sie die Gewinnschwelle (Break-even-Menge). \hfill (4~P.)
  \item \textbf{Entscheiden und begründen} Sie, ob das Projekt unter Kosten- und Erfolgsgesichtspunkten realisiert werden soll. \hfill (2~P.)
\end{enumerate}

\newpage
\section{Lösungen}\label{loesungen}

\begin{loesung}[title={Lösung Aufgabe~1 --- Zielsystem des Unternehmens}]
\textbf{(a)} \emph{Sachziele} beschreiben das konkrete Leistungsprogramm --- welche Güter
und Dienstleistungen in welcher Art und Menge erstellt werden. \emph{Formalziele}
(Erfolgsziele) sind die Maßstäbe, an denen der Erfolg gemessen wird (Gewinn,
Rentabilität, Wirtschaftlichkeit, Produktivität, Liquidität). Aussage (1) „hochwertige
Küchengeräte fertigen`` ist ein \emph{Sachziel}; Aussage (2) „Umsatzrendite auf 8~\%
steigern`` ist ein \emph{Formalziel}.

\textbf{(b)} Komplementäre (sich fördernde), konkurrierende (sich behindernde) und
indifferente (neutrale) Zielbeziehungen.

\textbf{(c)} Ziel (3) und Ziel (2): \emph{komplementär} --- eine halbierte Ausschussrate
senkt die Kosten und erhöht damit tendenziell die Umsatzrendite. Ziel (4) und Ziel (2):
kurzfristig eher \emph{konkurrierend}, weil Investitionen in emissionsärmere Fertigung
zunächst Kosten verursachen; langfristig kann die Beziehung komplementär werden
(Energieeinsparung, Image).
\end{loesung}

\begin{loesung}[title={Lösung Aufgabe~2 --- Produktionsfaktoren nach Gutenberg}]
\textbf{(a)} Die \emph{Volkswirtschaftslehre} betrachtet die gesamte Volkswirtschaft und
unterscheidet die Faktoren \emph{Arbeit, Boden und Kapital} (zunehmend ergänzt um
Wissen). Die \emph{Betriebswirtschaftslehre} betrachtet den einzelnen Betrieb; nach
Gutenberg werden \emph{Elementarfaktoren} (ausführende Arbeit, Betriebsmittel,
Werkstoffe) und der \emph{dispositive Faktor} (leitende Arbeit) unterschieden. Der
volkswirtschaftliche Faktor Arbeit wird also in ausführende und leitende (dispositive)
Arbeit aufgespalten; Boden und Kapital gehen in Betriebsmitteln und Werkstoffen auf.

\textbf{(b)} Elementarfaktoren: ausführende (objektbezogene) Arbeit, Betriebsmittel,
Werkstoffe. Dispositiver Faktor: Geschäfts- und Betriebsleitung mit den Aufgaben
Planung, Organisation und Kontrolle --- er kombiniert die Elementarfaktoren
zielgerichtet.

\textbf{(c)} Holz = \emph{Rohstoff} (Hauptbestandteil des Endprodukts); Leim =
\emph{Hilfsstoff} (geht ein, aber nur untergeordnet); Schmieröl = \emph{Betriebsstoff}
(wird verbraucht, geht nicht in das Produkt ein); CNC-Fräse = \emph{Betriebsmittel}
(wird gebraucht, nicht verbraucht).
\end{loesung}

\begin{loesung}[title={Lösung Aufgabe~3 --- Fertigungstypen}]
\textbf{(a)} \emph{Einzelfertigung} --- ein Unikat wird individuell nach Kundenauftrag
gefertigt.
\textbf{(b)} \emph{Sortenfertigung} --- artgleiche Varianten (Biersorten) aus demselben
Ausgangsmaterial auf derselben Anlage.
\textbf{(c)} \emph{Serienfertigung} --- begrenzte Stückzahl artverschiedener Produkte
(Modelle) nacheinander mit Umrüstung.
\textbf{(d)} \emph{Massenfertigung} --- sehr große Stückzahl eines gleichen Produkts
über langen Zeitraum.
\textbf{(e)} Bei der \emph{Sortenfertigung} sind die Erzeugnisse artgleich (Varianten
eines Grundprodukts aus gleichem Ausgangsmaterial), bei der \emph{Serienfertigung}
artverschieden (unterschiedliche Produkte auf denselben Anlagen).
\end{loesung}

\begin{loesung}[title={Lösung Aufgabe~4 --- Produktivität und Wirtschaftlichkeit}]
\textbf{(a)} Die \emph{Produktivität} ist ein Mengenverhältnis: Ausbringungsmenge
(Output) je Einsatzmenge (Input), z.\,B. Stück je Arbeitsstunde. Die
\emph{Wirtschaftlichkeit} ist ein Wertverhältnis: Ertrag (in €) geteilt durch Aufwand
(in €); sie berücksichtigt also Preise und Kosten, nicht nur Mengen.

\textbf{(b)} Arbeitsstunden je Monat: $20 \cdot 160 = 3.200$ Stunden.
August: $8.000 / 3.200 = 2{,}5$ Stück je Stunde.
September: $8.800 / 3.200 = 2{,}75$ Stück je Stunde --- die Produktivität ist gestiegen.

\textbf{(c)} August: $400.000 / 320.000 = 1{,}25$. September: $440.000 / 352.000 =
1{,}25$. Beide Monate sind wirtschaftlich (Wert $> 1$: je eingesetztem Euro werden
1,25~€ erwirtschaftet). Trotz gestiegener Produktivität blieb die Wirtschaftlichkeit
unverändert, weil Ertrag und Aufwand im gleichen Verhältnis gestiegen sind.

\textbf{(d)} Die \emph{Qualität} (Erfüllung der Kundenanforderungen, geringe
Ausschussraten).
\end{loesung}

\begin{loesung}[title={Lösung Aufgabe~5 --- Logistik: Teilbereiche und 6 R}]
\textbf{(a)} Die \emph{richtigen Güter} in der \emph{richtigen Menge} im
\emph{richtigen Zustand} zum \emph{richtigen Zeitpunkt} am \emph{richtigen Ort} zu den
\emph{richtigen Kosten}.

\textbf{(b)} Richtig ist: Die Beschaffungslogistik beschäftigt sich mit dem
Informations- und Materialfluss \emph{vom Lieferanten zum Unternehmen}. Die Abläufe im
Lager betreffen die \emph{Lagerlogistik}, die Abläufe in der Produktion die
\emph{Produktionslogistik}.

\textbf{(c)} Lieferant $\rightarrow$ (1) Materiallager $\rightarrow$ (2) Produktion
$\rightarrow$ (3) Absatzlager $\rightarrow$ (4) Versand $\rightarrow$ Kunde.
\end{loesung}

\begin{loesung}[title={Lösung Aufgabe~6 --- Lagerfunktionen und SCM}]
\textbf{(a)} \emph{Sicherungsfunktion (Vorratsfunktion):} Absicherung gegen
Lieferverzögerungen und Bedarfsschwankungen (z.\,B. Sicherheitsbestand an Rohstoffen).
\emph{Ausgleichsfunktion (Pufferfunktion):} Ausgleich unterschiedlicher Rhythmen von
Beschaffung, Produktion und Absatz (z.\,B. saisonale Ernte --- ganzjähriger Verbrauch).
\emph{Spekulationsfunktion:} Einkauf größerer Mengen bei erwarteten Preissteigerungen
oder zur Nutzung von Mengenrabatten. \emph{Veredelungsfunktion (Umformungsfunktion):}
Lagerung als Teil des Produktionsprozesses (Reifung von Käse oder Wein, Trocknung von
Holz). (Alternativ: \emph{Sortimentsfunktion} --- Zusammenstellung eines
verkaufsfähigen Sortiments im Handel.)

\textbf{(b)} Supply Chain Management ist die integrierte Planung und Steuerung der
gesamten Liefer- und Wertschöpfungskette --- vom Lieferanten über Hersteller,
Logistikdienstleister, Zwischenhändler und Wiederverkäufer bis zum Endkunden ---, bei
der Güter-, Informations- und Geldflüsse aller Beteiligten aufeinander abgestimmt
werden. Ziele (zwei genügen): kürzere Lieferzeiten, geringere Bestände entlang der
Kette, höhere Liefertreue, Vermeidung aufschaukelnder Bestellmengen
(Peitscheneffekt), Kostenvorteile für alle Partner.
\end{loesung}

\begin{loesung}[title={Lösung Aufgabe~7 --- Marketingbegriff, Marktforschung, Ziele}]
\textbf{(a)} Die Aussage ist \emph{falsch}. Ausgangspunkt aller Marketingaktivitäten
sind \emph{Informationen über den Markt} --- über Zielgruppe, Bedürfnisse und
Absatzwege. Das moderne Marketingdenken stellt den \emph{Kunden} und dessen
Problemlösung in den Mittelpunkt, nicht das Produkt.

\textbf{(b)} \emph{Primärforschung} (field research): Erhebung neuer, eigener Daten ---
Methoden: Befragung, Beobachtung, Experiment/Markttest, Panel.
\emph{Sekundärforschung} (desk research): Auswertung bereits vorhandener Daten ---
Quellen: interne Absatzstatistiken, amtliche Statistik, Verbands- und
Marktforschungsstudien. Sekundärforschung ist schneller und günstiger, aber oft nicht
passgenau und ggf. veraltet.

\textbf{(c)} \emph{Ökonomische Ziele} sind messbare wirtschaftliche Größen ---
Beispiele: Umsatz, Absatz, Gewinn, Marktanteil. \emph{Psychografische Ziele} sind nicht
direkt messbare, beim Kunden verankerte Größen --- Beispiele: Bekanntheitsgrad, Image,
Kundenzufriedenheit, Kundenbindung.
\end{loesung}

\begin{loesung}[title={Lösung Aufgabe~8 --- Fallaufgabe Marketing-Mix}]
\textbf{(a)} Produkt- und Leistungspolitik (Product), Preis- und Kontrahierungspolitik
(Price), Distributionspolitik (Place), Kommunikationspolitik (Promotion).

\textbf{(b)} Beispiele (je zwei Maßnahmen): \emph{Produktpolitik:} hochwertiges Design
und langlebiger Motor; Garantieverlängerung und Rezeptheft/App als Serviceleistung.
\emph{Preispolitik:} Premiumpreis passend zur Präferenzstrategie (Zielgruppe mit
höherem Einkommen); Einführungsrabatt oder Bundle-Angebote, Skonto für schnelle
Zahlung. \emph{Distributionspolitik:} Verkauf über den eigenen Onlineshop (direkter
Absatz), ergänzend ausgewählte Elektrofachmärkte (indirekter Absatz) zur Erhöhung des
Distributionsgrads. \emph{Kommunikationspolitik:} Social-Media-Kampagne mit
Food-Influencern; Verkaufsförderung durch Vorführaktionen und Produkttests am Point of
Sale. (Andere sachgerechte Maßnahmen sind gleichwertig.)

\textbf{(c)} \emph{Distributionspolitik} --- sie umfasst neben der Wahl der Absatzwege
die physische Distribution (Marketinglogistik), also die Aufgabe, das Produkt zum
richtigen Zeitpunkt in der gewünschten Menge an den Ort der Nachfrage zu bringen.

\textbf{(d)} \emph{Werbung} ist bezahlte, unmittelbar produkt- bzw. absatzbezogene
Kommunikation (Ziel: Kaufauslösung). \emph{Public Relations} (Öffentlichkeitsarbeit)
richtet sich auf das gesamte Unternehmen und will Vertrauen und ein positives Image in
der Öffentlichkeit aufbauen --- ohne unmittelbaren Produktbezug.
\end{loesung}

\begin{loesung}[title={Lösung Aufgabe~9 --- Erfolgsrechnung zur Produkteinführung}]
\textbf{(a)} Deckungsbeitrag je Stück: $30 - 14 = 16$~€. Gesamtdeckungsbeitrag:
$150.000 \cdot 16 = 2.400.000$~€. Gewinn: $2.400.000 - 150.000 = 2.250.000$~€.

\textbf{(b)} Break-even-Menge: Fixkosten geteilt durch Stückdeckungsbeitrag $=
150.000 / 16 = 9.375$ Stück.

\textbf{(c)} Das Projekt soll realisiert werden: Die Gewinnschwelle von 9.375 Stück
liegt weit unter der erwarteten Absatzmenge von 150.000 Stück; es wird ein deutlicher
Gewinn von 2,25~Mio.~€ erwartet. Selbst bei erheblich geringerem Absatz bliebe das
Projekt profitabel.
\end{loesung}

\vspace{1em}
\hrule
\vspace{0.8em}
\noindent\textsf{\small Weiterüben: Übungsband Betriebswirtschaft Kap.~1 (mit Lösungsteil) und das
interaktive Quiz dieser Lektion. Nächste Lektion: Betriebliche Funktionen II --- Rechnungswesen,
Finanzierung, Controlling \& Personal.}

\end{document}
